\chapter{Discussion}
\label{ch:discussion}

\section{Synthetic benchmarks can mislead}
% - Same code scores 0.9317 on our synthetic data and 0.5482 on LOFAR.
% - The synthetic label is noise-free (> 0.5 sigma by construction); real
%   labels are noisy and a fifth of real RFI has no brightness signal.
% - Components that look reasonable on paper did nothing on real data, and
%   the earlier synthetic ablation was too noisy to show it either.

\section{Lessons about measurement}
% Three times in this project a single run overstated an effect that
% shrank or vanished with more seeds:
%   lr 1e-4 gain  +0.118 (1 seed) -> +0.058, not significant (3 seeds) [PART 12.10]
%   GroupNorm     -0.0335 (1 seed) -> -0.009, p = 0.48 (3 seeds)       [PART 20]
%   strips on faint RFI: ~+0.02 at own thresholds -> 0 at matched budget [PART 18.5]
% Also the tf_unet failure itself was first blamed on the architecture.
% This section is about method; keep it short and factual.

\section{Limitations}
% [PART 14.8] -- state all of these plainly:
%  1. Published baselines (RFI-Net, their U-Net) are QUOTED from Mesarcik
%     et al. Table 2, not re-run. Our AOFlagger number matches theirs to
%     four decimals, so the test data and metric are the same -- but write
%     "as reported by Mesarcik et al. (2022)".
%  2. 109 test images: +/-0.04 bootstrap interval. Biggest uncertainty.
%  3. Training budgets differ from theirs (they: 32x32 patches, 100 epochs;
%     us: full images, 28,000 steps).
%  4. 3 seeds is the minimum; several conclusions rest on it.
%  5. tf_unet has not been run on HERA.
%  6. The gain over tf_unet is not yet attributed (Chapter 6).

\section{Future work}
% Depends on the novelty discussion. Measured candidates:
%  - Isolate the Dice loss (one run, ~25 min) [PART 19.2].
%  - Reduce false positives -- the largest measured headroom (+0.114).
%  - Apply to GMRT pulsar data (the original motivation).
